\subsection{Subjecthood as a Legal Category, Not a Scientific One}
\label{sec:9.3.1}
Section~\ref{sec:2.3.1} argued that sentience is not a threshold a system crosses but a set of separate, gradable capacities — nociception, learned avoidance, a self-model, something like suffering — each with its own uncertain evidence. That conclusion is a problem for a legal system specifically, because law generally needs a decidable rule: a court or a regulator has to say whether a given system is or is not a covered subject, on a specific date, even when the underlying science offers a spectrum with no natural cut point.

Law has faced this exact mismatch before and did not solve it by waiting for the science to resolve. Fetal viability, brain death, and the age of criminal responsibility are all legal thresholds drawn across what is actually a continuous biological process, defended not as scientific facts but as the least-bad place to put an enforceable line given what's at stake on either side of it. The same move is available here: a jurisdiction can set an operational threshold — a specific, testable benchmark on self-monitoring, learned aversive response, or behavioral flexibility under novel conditions — without claiming that benchmark settles what sentience actually is, only that it is where legal protection starts to attach.

Two designs are available, and they trade off differently. A single bright-line threshold is easy to enforce and easy to game: a system built to clear the specific test without having whatever the test was meant to detect. A graded framework — obligations that scale with how many benchmarks a system meets, echoing that section's own suggestion that protection should track capacity rather than a baseline bar — is harder to game but harder to adjudicate, since a court now has to weigh degree instead of answering yes or no.

Either design needs a body to apply it. Its composition is already specified: AI researchers, consciousness scientists, ethicists, and legal scholars, interdisciplinary by design because no single field's expertise is sufficient on its own. What that leaves open is what happens once the body rules — whether its judgment carries any force beyond its own credibility.
